\begin{bpnew}{propborder}{Definition: Sequence}
\begin{definition}[Sequence]
\label{def:sequence}

Let $X$ be a nonempty set. A \textbf{sequence} in $X$ is a function
\[
    x : \mathbb{N} \to X.
\]
The value $x(n)$ is written $x_n$ and called the \textbf{$n$-th term} of the
sequence. The sequence itself is denoted $(x_n)_{n \in \mathbb{N}}$, or
simply $(x_n)$ when the index set is clear.
\end{definition}
\end{bpnew}

\begin{remark*}[Standard quantified statement]
\[
\forall X \neq \varnothing \;\forall x : \mathbb{N} \to X
\;\bigl(x \text{ is a sequence in } X\bigr).
\]
\end{remark*}

\begin{remark*}[Definition predicate reading]
\[
\operatorname{Sequence}(x_n, X)
\;:=\;
x : \mathbb{N} \to X.
\]
\end{remark*}

\begin{remark*}[Interpretation]
A sequence is not a set of terms but a function: the index carries
independent weight. Two sequences $(x_n)$ and $(y_n)$ in $X$ are equal
precisely when $x_n = y_n$ for every $n \in \mathbb{N}$, so order and
repetition are intrinsic to the object. The notational collapse from
$x(n)$ to $x_n$ is conventional suppression of function-application
syntax; the underlying object remains a function $\mathbb{N} \to X$.
When $X = \mathbb{R}$, the sequence is called a \textbf{real sequence},
which is the default ambient context for this chapter.
\end{remark*}











\begin{example}[Constant Sequence]
\label{ex:constant-sequence}

Let $c \in \mathbb{R}$. The \textbf{constant sequence} is defined by
\[
    x_n := c \quad \forall n \in \mathbb{N}.
\]
The constant sequence converges to $c$: given any $\varepsilon > 0$,
take $N = 1$, and observe $|x_n - c| = 0 < \varepsilon$ for all $n \ge 1$.
It is the degenerate case of convergence in which the tail is not merely
eventually close to $L$ but identically equal to it from the first term.
\end{example}

\begin{example}[The Sequence $(1/n)$]
\label{ex:sequence-reciprocal}

The sequence $(x_n)$ defined by
\[
    x_n := \frac{1}{n} \quad \forall n \in \mathbb{N}
\]
converges to $0$. Given $\varepsilon > 0$, the Archimedean property
supplies $N \in \mathbb{N}$ with $N > 1/\varepsilon$, and then for all
$n \ge N$,
\[
    \left|\frac{1}{n} - 0\right| = \frac{1}{n} \le \frac{1}{N} < \varepsilon.
\]
This is the canonical example of a sequence that approaches its limit
monotonically from above, never reaching it. It also serves as the
standard witness in squeeze arguments and in the proof that
$\inf\{1/n : n \in \mathbb{N}\} = 0$.
\end{example}

\begin{example}[The Sequence $(n)$]
\label{ex:sequence-natural}

The sequence $(x_n)$ defined by
\[
    x_n := n \quad \forall n \in \mathbb{N}
\]
is divergent. For any proposed limit $L \in \mathbb{R}$, take
$\varepsilon = 1$; the Archimedean property supplies $n \in \mathbb{N}$
with $n > L + 1$, so $|x_n - L| = |n - L| > 1 = \varepsilon$ for
infinitely many $n$. The sequence is unbounded, and every convergent
sequence is bounded, so divergence also follows from that theorem once
it is established. This is the canonical example of divergence to
$+\infty$: the terms are not oscillating outside a neighborhood of $L$
but growing without bound in a single direction.
\end{example}




































\begin{bpnew}{propborder}{Definition: Convergence of a Sequence}
\begin{definition}[Convergence of a Sequence]
\label{def:convergent-sequence}

Let $(x_n) \subseteq \mathbb{R}$. The sequence $(x_n)$ \textbf{converges}
to $L \in \mathbb{R}$ if
\[
    \forall \varepsilon > 0 \;\exists N \in \mathbb{N} \;\forall n \ge N
    \;\bigl(|x_n - L| < \varepsilon\bigr).
\]
The value $L$ is called the \textbf{limit} of $(x_n)$, written
$x_n \to L$ or $\displaystyle\lim_{n \to \infty} x_n = L$.
A sequence that converges to some $L \in \mathbb{R}$ is called
\textbf{convergent}; otherwise it is called \textbf{divergent}.
\end{definition}
\end{bpnew}

\begin{remark*}[Standard quantified statement]
\[
\begin{aligned}
&\forall (x_n) \subseteq \mathbb{R} \;\forall L \in \mathbb{R}\\
&\quad\Bigl(
    x_n \to L
    \;\Longleftrightarrow\;
    \forall \varepsilon > 0 \;\exists N \in \mathbb{N} \;\forall n \ge N
    \;\bigl(|x_n - L| < \varepsilon\bigr)
\Bigr).
\end{aligned}
\]
\end{remark*}

\begin{remark*}[Definition predicate reading]
\[
\operatorname{ConvergentSequence}(x_n, L)
\;:=\;
\forall \varepsilon > 0 \;\exists N \in \mathbb{N} \;\forall n \ge N
\;\operatorname{OutputTolerance}(x_n, L, \varepsilon).
\]
\end{remark*}

\begin{remark*}[Negated quantified statement]
\[
\begin{aligned}
x_n \not\to L
\;\Longleftrightarrow\;
&\;\exists \varepsilon > 0 \;\forall N \in \mathbb{N} \;\exists n \ge N\\
&\quad\bigl(|x_n - L| \ge \varepsilon\bigr).
\end{aligned}
\]
\end{remark*}

\begin{remark*}[Negation predicate reading]
\[
\operatorname{DivergentSequence}(x_n)
\;:=\;
\exists \varepsilon > 0 \;\forall N \in \mathbb{N} \;\exists n \ge N
\;\bigl(|x_n - L| \ge \varepsilon\bigr).
\]
\end{remark*}

\begin{remark*}[Failure mode decomposition]
The negation decomposes into a single structural obstruction:
\[
\underbrace{\exists \varepsilon > 0}_{\text{a fixed tolerance}}
\;\underbrace{\forall N \in \mathbb{N}}_{\text{no tail suffices}}
\;\underbrace{\exists n \ge N}_{\text{terms escape infinitely often}}
\;\underbrace{|x_n - L| \ge \varepsilon}_{\text{outside the target interval}}.
\]
The sequence fails to converge to $L$ not by drifting away once, but
by returning outside the $\varepsilon$-interval infinitely often regardless
of how far along the tail one looks.
\end{remark*}

\begin{remark*}[Interpretation]
The quantifier order $\forall \varepsilon \;\exists N \;\forall n \ge N$
encodes a strict dependency: $N$ is permitted to depend on $\varepsilon$,
but not on $n$. This asymmetry is load-bearing. Reversing the outer
quantifiers to $\exists N \;\forall \varepsilon$ would assert that a
single index $N$ controls all tolerances simultaneously, which is a
strictly stronger and generally false condition. The condition
$|x_n - L| < \varepsilon$ is the membership assertion $x_n \in (L -
\varepsilon,\, L + \varepsilon)$; convergence is therefore the statement
that the tail of the sequence is eventually contained in every
$\varepsilon$-neighborhood of $L$.
\end{remark*}

\begin{bpnew}{propborder}{Definition: Convergence --- Neighborhood Formulation}
\begin{definition}[Convergence --- Neighborhood Formulation]
\label{def:convergent-sequence-neighborhood}

Let $(x_n) \subseteq \mathbb{R}$ and let $L \in \mathbb{R}$. The sequence
$(x_n)$ \textbf{converges} to $L$ if for every $\varepsilon > 0$ the
$\varepsilon$-neighborhood
\[
    V_\varepsilon(L) \;:=\; (L - \varepsilon,\, L + \varepsilon)
\]
contains all but finitely many terms of $(x_n)$; that is,
\[
    \forall \varepsilon > 0 \;\exists N \in \mathbb{N} \;\forall n \ge N
    \;\bigl(x_n \in V_\varepsilon(L)\bigr).
\]
\end{definition}
\end{bpnew}

\begin{remark*}[Standard quantified statement]
\[
\begin{aligned}
&\forall (x_n) \subseteq \mathbb{R} \;\forall L \in \mathbb{R}\\
&\quad\Bigl(
    x_n \to L
    \;\Longleftrightarrow\;
    \forall \varepsilon > 0 \;\exists N \in \mathbb{N} \;\forall n \ge N
    \;\bigl(x_n \in V_\varepsilon(L)\bigr)
\Bigr).
\end{aligned}
\]
\end{remark*}

\begin{remark*}[Definition predicate reading]
\[
\operatorname{ConvergentSequence}(x_n, L)
\;:=\;
\forall \varepsilon > 0 \;\exists N \in \mathbb{N} \;\forall n \ge N
\;\bigl(x_n \in V_\varepsilon(L)\bigr).
\]
\end{remark*}

\begin{remark*}[Negated quantified statement]
\[
\begin{aligned}
x_n \not\to L
\;\Longleftrightarrow\;
&\;\exists \varepsilon > 0 \;\forall N \in \mathbb{N} \;\exists n \ge N\\
&\quad\bigl(x_n \notin V_\varepsilon(L)\bigr).
\end{aligned}
\]
\end{remark*}

\begin{remark*}[Negation predicate reading]
\[
\operatorname{DivergentSequence}(x_n)
\;:=\;
\exists \varepsilon > 0 \;\forall N \in \mathbb{N} \;\exists n \ge N
\;\bigl(x_n \notin V_\varepsilon(L)\bigr).
\]
\end{remark*}

\begin{remark*}[Failure mode decomposition]
\[
\underbrace{\exists \varepsilon > 0}_{\text{a fixed neighborhood}}
\;\underbrace{\forall N \in \mathbb{N}}_{\text{no tail suffices}}
\;\underbrace{\exists n \ge N}_{\text{terms escape infinitely often}}
\;\underbrace{x_n \notin V_\varepsilon(L)}_{\text{outside every candidate neighborhood}}.
\]
The sequence fails to converge to $L$ precisely when some neighborhood
of $L$ contains only finitely many terms: infinitely many terms lie
permanently outside it.
\end{remark*}

\begin{remark*}[Interpretation]
The neighborhood formulation reframes convergence as a set-containment
condition rather than an inequality condition. The assertion $x_n \in
V_\varepsilon(L)$ is identical to $|x_n - L| < \varepsilon$; the two
definitions are logically equivalent and the choice between them is one
of language, not content. The neighborhood language is the natural
bridge to the topological formulation: in a general topological space,
open sets replace $\varepsilon$-neighborhoods and the definition reads
identically with $V_\varepsilon(L)$ replaced by any open set $U$
containing $L$. The phrase \emph{all but finitely many terms} is the
canonical prose compression of the quantifier block $\exists N \in
\mathbb{N} \;\forall n \ge N$, and is standard in the literature.
\end{remark*}










\begin{bpnew}{thmborder}{Theorem: Equivalence of Convergence Formulations}
\begin{theorem}[Equivalence of Convergence Formulations]
\label{thm:equivalence-of-convergence-formulations}

Let $(x_n) \subseteq \mathbb{R}$ and let $L \in \mathbb{R}$. The following
statements are equivalent.
\begin{enumerate}[label=\textup{(\alph*)}]
    \item $x_n \to L$.
    \item $\forall \varepsilon > 0 \;\exists K \in \mathbb{N} \;\forall n \ge K
          \;\bigl(|x_n - L| < \varepsilon\bigr)$.
    \item $\forall \varepsilon > 0 \;\exists K \in \mathbb{N} \;\forall n \ge K
          \;\bigl(L - \varepsilon < x_n < L + \varepsilon\bigr)$.
    \item $\forall \varepsilon > 0 \;\exists K \in \mathbb{N} \;\forall n \ge K
          \;\bigl(x_n \in V_\varepsilon(L)\bigr)$.
\end{enumerate}

\smallskip
\noindent
\hyperref[prf:equivalence-of-convergence-formulations]{\textit{Go to proof.}}
\end{theorem}
\end{bpnew}

\begin{remark*}[Standard quantified statement]
\[
\begin{aligned}
&\forall (x_n) \subseteq \mathbb{R} \;\forall L \in \mathbb{R}\\
&\quad\Bigl(
    (a) \;\Longleftrightarrow\; (b)
         \;\Longleftrightarrow\; (c)
         \;\Longleftrightarrow\; (d)
\Bigr).
\end{aligned}
\]
\end{remark*}

\begin{remark*}[Interpretation]
The four conditions are the same geometric fact written in four
notational registers. Condition (a) is the informal assertion.
Condition (b) is the canonical $\varepsilon$-$K$ definition: the
absolute value formulation. Condition (c) makes the interval
structure explicit by expanding $|x_n - L| < \varepsilon$ into its
equivalent double inequality $L - \varepsilon < x_n < L +
\varepsilon$; this is useful in proofs where upper and lower bounds
are handled separately. Condition (d) is the neighborhood
formulation: $x_n \in V_\varepsilon(L)$ is identical to conditions
(b) and (c) but frames convergence as set membership, which is the
natural language for the topological generalization. The index $K$
plays the same role as $N$ in the base definition
\hyperref[def:convergent-sequence]{Definition~\ref*{def:convergent-sequence}};
the letter changes here to signal that $K$ is the tail offset in
the sense of \hyperref[def:m-tail]{Definition~\ref*{def:m-tail}}.
\end{remark*}






















\begin{bpnew}{thmborder}{Theorem: Uniqueness of Limits}
\begin{theorem}[Uniqueness of Limits]
\label{thm:uniqueness-of-limits}

Let $(x_n) \subseteq \mathbb{R}$. If $x_n \to L$ and $x_n \to K$, then
$L = K$.

\smallskip
\noindent
\hyperref[prf:uniqueness-of-limits]{\textit{Go to proof.}}
\end{theorem}
\end{bpnew}

\begin{remark*}[Standard quantified statement]
\[
\begin{aligned}
&\forall (x_n) \subseteq \mathbb{R} \;\forall L \in \mathbb{R}
\;\forall K \in \mathbb{R}\\
&\quad\Bigl(
    \bigl(x_n \to L \;\land\; x_n \to K\bigr)
    \;\Longrightarrow\;
    L = K
\Bigr).
\end{aligned}
\]
\end{remark*}

\begin{remark*}[Negated quantified statement]
\[
\begin{aligned}
&\exists (x_n) \subseteq \mathbb{R} \;\exists L \in \mathbb{R}
\;\exists K \in \mathbb{R}\\
&\quad\Bigl(
    \bigl(x_n \to L \;\land\; x_n \to K\bigr)
    \;\land\;
    L \neq K
\Bigr).
\end{aligned}
\]
\end{remark*}

\begin{remark*}[Failure mode decomposition]
\[
\underbrace{\exists (x_n) \subseteq \mathbb{R}}_{\text{some sequence}}
\;\underbrace{x_n \to L \;\land\; x_n \to K}_{\text{converges to two values}}
\;\underbrace{L \neq K}_{\text{distinct limits}}.
\]
The theorem asserts this configuration is impossible in $\mathbb{R}$.
The failure mode is not merely exotic: it occurs in non-Hausdorff
topological spaces where distinct points cannot be separated by disjoint
open sets. The proof that it cannot occur in $\mathbb{R}$ is precisely
an argument that $\mathbb{R}$ has enough room to separate any two
distinct points by disjoint $\varepsilon$-neighborhoods.
\end{remark*}

\begin{remark*}[Contrapositive quantified statement]
\[
\begin{aligned}
&\forall (x_n) \subseteq \mathbb{R} \;\forall L \in \mathbb{R}
\;\forall K \in \mathbb{R}\\
&\quad\Bigl(
    L \neq K
    \;\Longrightarrow\;
    \lnot\bigl(x_n \to L \;\land\; x_n \to K\bigr)
\Bigr).
\end{aligned}
\]
\end{remark*}

\begin{remark*}[Interpretation]
The theorem licenses the notation $\displaystyle\lim_{n \to \infty} x_n
= L$: the definite article is justified precisely because the limit,
when it exists, is unique. Without uniqueness the limit would be a
relation rather than a function, and the notation would be ambiguous.
The proof strategy is a separation argument — if $L \neq K$ then
$|L - K| > 0$, and choosing $\varepsilon = |L - K|/2$ produces two
disjoint neighborhoods that the tail of $(x_n)$ cannot simultaneously
occupy. The Hausdorff remark in the failure mode block is the forward
pointer: when sequences are studied in general topological spaces,
uniqueness of limits is equivalent to the Hausdorff separation axiom.
\end{remark*}





\begin{bpnew}{propborder}{Definition: $M$-Tail of a Sequence}
\begin{definition}[$M$-Tail of a Sequence]
\label{def:m-tail}

Let $(x_n) \subseteq \mathbb{R}$ and let $M \in \mathbb{N}$. The
\textbf{$M$-tail} of $(x_n)$ is the sequence
\[
    (x_n)_{n \ge M} \;:=\; \bigl(x_M,\, x_{M+1},\, x_{M+2},\, \ldots\bigr).
\]
\end{definition}
\end{bpnew}

\begin{remark*}[Standard quantified statement]
\[
\forall (x_n) \subseteq \mathbb{R} \;\forall M \in \mathbb{N}
\;\Bigl(
    (x_n)_{n \ge M} = \{x_n : n \ge M\}
\Bigr).
\]
\end{remark*}

\begin{remark*}[Interpretation]
The $M$-tail discards the first $M - 1$ terms and retains everything
from index $M$ onward. It is the formal object underlying the phrase
\emph{all but finitely many terms}: to say that all but finitely many
terms of $(x_n)$ satisfy a property $P$ is precisely to say that some
$M$-tail of $(x_n)$ satisfies $P$ termwise. The index $M$ is an offset
into the sequence, not a bound on the terms themselves; two sequences
with identical $M$-tails share all asymptotic properties, differing
only in finitely many initial values which convergence arguments
discard entirely.
\end{remark*}



\begin{bpnew}{thmborder}{Theorem: Convergence of a Tail}
\begin{theorem}[Convergence of a Tail]
\label{thm:convergence-of-tail}

Let $(x_n) \subseteq \mathbb{R}$ and let $m \in \mathbb{N}$. The
$m$-tail $(x_{m+n})_{n \in \mathbb{N}}$ converges if and only if
$(x_n)$ converges. Moreover, in this case,
\[
    \lim_{n \to \infty} x_{m+n} \;=\; \lim_{n \to \infty} x_n.
\]

\smallskip
\noindent
\hyperref[prf:convergence-of-tail]{\textit{Go to proof.}}
\end{theorem}
\end{bpnew}

\begin{remark*}[Standard quantified statement]
\[
\begin{aligned}
&\forall (x_n) \subseteq \mathbb{R} \;\forall m \in \mathbb{N}\\
&\quad\Bigl(
    \bigl(x_{m+n}\bigr)_{n \in \mathbb{N}} \text{ converges}
    \;\Longleftrightarrow\;
    (x_n) \text{ converges}
\Bigr)\\
&\quad\land\;\Bigl(
    \lim_{n\to\infty} x_{m+n} = \lim_{n\to\infty} x_n
\Bigr).
\end{aligned}
\]
\end{remark*}

\begin{remark*}[Contrapositive quantified statement]
\[
\begin{aligned}
&\forall (x_n) \subseteq \mathbb{R} \;\forall m \in \mathbb{N}\\
&\quad\Bigl(
    \bigl(x_{m+n}\bigr)_{n \in \mathbb{N}} \text{ diverges}
    \;\Longleftrightarrow\;
    (x_n) \text{ diverges}
\Bigr).
\end{aligned}
\]
\end{remark*}

\begin{remark*}[Interpretation]
Convergence is a \emph{tail property}: it depends only on the
eventual behavior of the sequence, not on any finite initial
segment. Discarding the first $m$ terms neither creates nor
destroys convergence, and when the limit exists it is preserved
exactly. This theorem formalizes the informal principle that
finitely many terms are irrelevant to asymptotic behavior. It is
the rigorous justification for the common proof move of assuming
without loss of generality that $n$ is large enough for some
auxiliary condition to hold — one simply passes to a suitable
$m$-tail, applies the argument there, and concludes for the full
sequence by this theorem. The contrapositive is equally useful:
to show $(x_n)$ diverges it suffices to show that some $m$-tail
diverges.
\end{remark*}









\begin{bpnew}{thmborder}{Theorem: Convergence by Domination}
\begin{theorem}[Convergence by Domination]
\label{thm:convergence-by-domination}

Let $(x_n) \subseteq \mathbb{R}$ and let $L \in \mathbb{R}$. Let
$(a_n)$ be a sequence of positive real numbers with
$\displaystyle\lim_{n \to \infty} a_n = 0$. If there exist a constant
$c > 0$ and an index $m \in \mathbb{N}$ such that
\[
    |x_n - L| \;\le\; c\, a_n \qquad \forall n \ge m,
\]
then $\displaystyle\lim_{n \to \infty} x_n = L$.

\smallskip
\noindent
\hyperref[prf:convergence-by-domination]{\textit{Go to proof.}}
\end{theorem}
\end{bpnew}

\begin{remark*}[Standard quantified statement]
\[
\begin{aligned}
&\forall (x_n) \subseteq \mathbb{R} \;\forall L \in \mathbb{R}
 \;\forall (a_n) \subseteq (0, \infty) \;\forall c > 0
 \;\forall m \in \mathbb{N}\\
&\quad\Bigl(
    \Bigl(
        \lim_{n \to \infty} a_n = 0
        \;\land\;
        \forall n \ge m \;\bigl(|x_n - L| \le c\, a_n\bigr)
    \Bigr)
    \;\Longrightarrow\;
    \lim_{n \to \infty} x_n = L
\Bigr).
\end{aligned}
\]
\end{remark*}

\begin{remark*}[Contrapositive quantified statement]
\[
\begin{aligned}
&\forall (x_n) \subseteq \mathbb{R} \;\forall L \in \mathbb{R}
 \;\forall (a_n) \subseteq (0,\infty) \;\forall c > 0
 \;\forall m \in \mathbb{N}\\
&\quad\Bigl(
    \lim_{n \to \infty} x_n \neq L
    \;\Longrightarrow\;
    \lim_{n \to \infty} a_n \neq 0
    \;\lor\;
    \exists n \ge m \;\bigl(|x_n - L| > c\, a_n\bigr)
\Bigr).
\end{aligned}
\]
\end{remark*}

\begin{remark*}[Interpretation]
The theorem is a \emph{domination principle}: it is not necessary to
verify the $\varepsilon$-$N$ condition directly for $(x_n)$. It
suffices to find a null sequence $(a_n)$ — one converging to zero —
that dominates the error $|x_n - L|$ on some tail, up to a fixed
scalar $c$. The scalar $c$ is irrelevant to the conclusion: since
$a_n \to 0$, the product $c\, a_n \to 0$ regardless of the size of
$c$, and any $\varepsilon > 0$ is eventually beaten by $c\, a_n$.
This theorem is the rigorous engine behind the common proof pattern
of bounding $|x_n - L|$ by a simpler expression and observing that
the simpler expression tends to zero. The canonical instance is
$a_n = 1/n$, where domination by $c/n$ is sufficient to conclude
convergence to $L$.
\end{remark*}

































